% Run pdflatex until references stabilize. The bibliography is part of this file.
\documentclass[11pt,a4paper]{article}
\usepackage[T1]{fontenc}
\usepackage{lmodern}
\usepackage{amsmath,amssymb,amsthm}
\usepackage[a4paper,left=25mm,right=25mm,top=19mm,bottom=20mm]{geometry}
\usepackage{xcolor}
\definecolor{linkink}{RGB}{18,70,78}
\usepackage[colorlinks=true,linkcolor=linkink,citecolor=linkink,urlcolor=linkink,
  pdftitle={Every finite joint measurability structure is realizable on one qubit},
  pdfauthor={Erfan Oliaei},pdfsubject={Quantum measurement incompatibility}]{hyperref}
\newtheorem{theorem}{Theorem}
\newtheorem{lemma}[theorem]{Lemma}
\newtheorem{corollary}[theorem]{Corollary}
\theoremstyle{definition}
\newtheorem*{example}{Example}
\newcommand{\Id}{I}
\newcommand{\C}{\mathbb C}
\newcommand{\R}{\mathbb R}
\newcommand{\tr}{\operatorname{tr}}
\newcommand{\op}{\mathrm{op}}
\newcommand{\K}{\mathcal K}
\newcommand{\D}{\mathcal D}
\setlength{\emergencystretch}{1.5em}
\title{\bfseries Every finite joint measurability structure\\ is realizable on one qubit}
\author{Erfan Oliaei}
\date{26 November 2026}

\begin{document}
\maketitle
\begin{abstract}
Every finite abstract simplicial complex is the family of jointly measurable
subsets of a collection of measurements on one qubit. This resolves a
conjecture of Andrejic and Kunjwal. For \(r\ge2\) maximal faces, the merged
construction uses at most \(r\) outcomes per measurement and gives real
rational effects in a common basis. Each measurement reports labels for
the maximal faces that omit its vertex. Distinct labels are mutually
exclusive in a joint measurement because their rank-one effects have
different ranges. The signed construction with suitably separated Bloch
axes uses at most \(2r-2\) outcomes per measurement and tolerates every
depolarizing noise fraction below \(1/(2r^2)\). We give explicit
bounds for independent perturbations of both constructions. At fixed
Hilbert space dimension and fixed outcome cap, only \(2^{O(n\log n)}\)
labelled pairwise compatibility graphs can occur on \(n\) measurements.
Thus no fixed outcome cap is universal in any fixed dimension.
\end{abstract}

\section{The question}
A joint measurement records one outcome for each observable in a single
experiment and reproduces each observable's statistics on every input state. Forgetting
coordinates gives a joint measurement for each subcollection. Jointly measurable subsets therefore form a simplicial
complex. Kunjwal, Heunen and Fritz proved the converse with binary
measurements on a Hilbert space whose dimension can grow with the
complex~\cite{KHF}. Andrejic and Kunjwal realized all structures on four
vertices and two infinite families on a qubit. They conjectured that
qubit measurements with arbitrary finite outcome alphabets can realize
every finite structure~\cite[Conjecture~4]{AK}. Kunjwal returned to this
question in 2025~\cite[Sec.~5.3]{HDR}. Section~\ref{sec:count} shows that
any fixed outcome cap excludes some structures.

A finite positive operator-valued measure (POVM) is a family \(M=(M(x))_{x\in X}\) of
positive semidefinite operators summing to the identity \(\Id\).
A labelled subfamily \((M_v)_{v\in S}\) is \emph{jointly measurable}, or
\emph{compatible}, if there is a POVM \(G\) on \(\prod_{v\in S}X_v\) such that
\[
     \sum_{\boldsymbol x:\,x_v=x}G(\boldsymbol x)=M_v(x)
     \qquad (v\in S,\ x\in X_v).
\]
We use the usual order \(A\preceq B\) for positive semidefinite \(B-A\).
A complex \(\K\) on a finite set \(V\) contains the empty set and all
singletons and is closed under subsets. Its maximal faces are its
\emph{facets}.

\begin{theorem}[Qubit universality]\label{thm:universal}
Every finite simplicial complex \(\K\) is the compatibility complex of a
family of POVMs on \(\C^2\). If \(\K\) has \(r\ge2\) facets, there is a
realization with at most \(r\) outcomes per POVM and with real rational
effects in a common basis. The least Hilbert space dimension is one
for a simplex and two for every other complex.
\end{theorem}

A measurement at \(v\) can report the labels of the facets that \emph{omit}
\(v\). A set of vertices is a face precisely when at least one label is
absent from their combined list. Assigning different spin bases to different
labels forces a parent measurement to count their weights separately.

\section{Encoding the facets}
Fix \(r\) orthonormal qubit bases \((p_{j,+},p_{j,-})\), with no ray shared
by distinct bases, and write \(P_{j,s}=p_{j,s}p_{j,s}^*\). Give each label
\(j\) a weight \(w_j>0\), and put \(w(A)=\sum_{j\in A}w_j\) for
\(A\subseteq[r]=\{1,\ldots,r\}\). Whenever \(w(A)\le1\), define
\begin{equation}\label{eq:measurement}
 M_A(j,s)=w_jP_{j,s}\quad (j\in A,\ s\in\{+,-\}),\qquad
 M_A(\bot)=(1-w(A))\Id .
\end{equation}
This is a POVM. With probability \(w_j\) it measures basis \(j\) and
reports both \(j\) and its sign. With the remaining probability it reports
\(\bot\). A zero outcome may always be omitted.

\begin{lemma}[Union criterion]\label{lem:union}
The measurements \(M_{A_1},\ldots,M_{A_\ell}\) are jointly measurable if and
only if \(w(A_1\cup\cdots\cup A_\ell)\le1\).
\end{lemma}
\begin{proof}
Write \(U=\bigcup_i A_i\). If \(w(U)\le1\), use \(M_U\) as parent. For
each \(i\), keep a label in \(A_i\) and send every other outcome to \(\bot\).

Conversely, suppose there is a joint POVM \(G\). For each \(j\in U\) choose
one index \(i(j)\) with \(j\in A_{i(j)}\). Let \(B_j\) be the event that
coordinate \(i(j)\) reports label \(j\) with either sign. Marginalizing gives
\(G(B_j)=w_j\Id\). For \(j\ne k\), the events are disjoint if \(i(j)=i(k)\)
because a coordinate reports one label at most. If \(i(j)\ne i(k)\), any
atom in \(B_j\cap B_k\) has an effect \(H\) bounded by both marginals,
\[
          0\preceq H\preceq w_jP_{j,s},\qquad H\preceq w_kP_{k,t},
\]
where \(s,t\) are the reported signs. A positive \(H\preceq wP\) vanishes
on \(\ker P\) because \(\|H^{1/2}x\|^2\le w\langle x,Px\rangle=0\) there.
Thus the range of \(H\) lies in both rays and \(H=0\). Each nonzero atom
contributes at most once, so \(\sum_{j\in U}G(B_j)=w(U)\Id\preceq\Id\).
\end{proof}

For singleton label sets with common weight, this criterion recovers the
\(1/r\) efficiency threshold for \(r\) lossy projective
measurements~\cite{SC}. Keeping the basis label also lets us group any
admissible collection of those tests into one observable.

\begin{proof}[Proof of Theorem~\ref{thm:universal}]
List the facets as \(F_1,\ldots,F_r\). For \(r\ge2\), set
\begin{equation}\label{eq:facetchoice}
       A_v=\{j:v\notin F_j\},\qquad w_j=c=\frac1{r-1},
       \qquad M_v=M_{A_v}.
\end{equation}
Each vertex belongs to a facet, so \(|A_v|\le r-1\). The union criterion
gives the two equivalences
\[
 (M_v)_{v\in S}\ \text{compatible}
 \quad\Longleftrightarrow\quad \bigcup_{v\in S}A_v\ne[r]
 \quad\Longleftrightarrow\quad S\subseteq F_j\ \text{for some }j .
\]
If \(|A_v|=r-1\) the erasure effect is zero, leaving \(2r-2\) outcomes.
Otherwise there are at most \(2r-3\). We call this the \emph{signed
realization}.

To reduce the alphabet to \(r\) outcomes, send every negative sign to
\(\bot\), leaving just one reported outcome per label in \(A_v\) and one
residual outcome. Explicitly, the reported effect is \(cP_j\) for each
\(j\in A_v\), and the residual effect is \(\Id-\sum_{j\in A_v}cP_j\).
Choose the positive rays close to \(e_1=(1,0)^{\mathsf T}\) by taking
\begin{equation}\label{eq:rational}
 P_j=P_{j,+}=\frac1{1+t_j^2}
       \begin{pmatrix}1&t_j\\ t_j&t_j^2\end{pmatrix},\qquad
 t_j=j/r^2,\qquad P_{j,-}=\Id-P_j .
\end{equation}
These bases share no ray, and their projectors have rational entries.
For \(U=\bigcup_{v\in S}A_v\), the merged measurements on \(S\) have a
parent exactly when
\begin{equation}\label{eq:mergedunion}
                    \sum_{j\in U}cP_j\preceq\Id .
\end{equation}
For necessity, take a joint POVM \(G\). For each \(j\in U\), choose a
vertex \(v_j\in S\) with \(j\in A_{v_j}\). Let \(B_j\) be the event that
\(v_j\) reports \(j\), so \(G(B_j)=cP_j\). The rank-one argument of
Lemma~\ref{lem:union} gives \(G(B_j\cap B_k)=0\) for \(j\ne k\), hence
\[
              \sum_{j\in U}cP_j=\sum_{j\in U}G(B_j)\preceq\Id.
\]
For sufficiency, the effects \(cP_j\), \(j\in U\), and their residual
form a parent. Each vertex keeps its own
labels and sends the rest to \(\bot\). Faces have \(|U|\le r-1\), so they
satisfy \eqref{eq:mergedunion}. A nonface has \(U=[r]\), and
\[
   \left\langle e_1,\Big(\sum_jcP_j\Big)e_1\right\rangle
       =c\sum_j\frac1{1+t_j^2}\ \ge\ \frac{rc}{1+r^{-2}}\ >\ 1 .
\]
This proves the bound of \(|A_v|+1\le r\) with real rational effects.

In dimension one the product of the individual distributions is a parent,
so the realizable complexes are exactly the simplices, including the one
on the empty vertex set. Scalar POVMs realize them on a qubit as well.
\end{proof}

\begin{example}
For facets \(\{1,2,3\}\), \(\{2,4\}\) and \(\{3,4\}\), the label sets
in vertex order are \(\{2,3\}\), \(\{3\}\), \(\{2\}\) and \(\{1\}\),
with weight \(1/2\) per label. Exactly two minimal sets display all three
labels, namely \(\{1,4\}\) and \(\{2,3,4\}\).
\end{example}

Cardinality thresholds already occur with binary measurements. For integers
\(1\le k<N\), give vertex \(j\in[N]\) the two effects \(P_j/k\) and
\(\Id-P_j/k\), using \eqref{eq:rational} with \(r=N\). The proof of
\eqref{eq:mergedunion} with weight \(1/k\) shows that \(S\) is compatible
exactly when \(\sum_{j\in S}P_j\preceq k\Id\). This holds for \(|S|\le k\).
For \(|S|\ge k+1\), the expectation of the sum in \(e_1\) is at least
\((k+1)/(1+N^{-2})>k\). Thus \(k=2\) gives, for every \(N\ge3\), the
\(N\)-complete scenario in which every pair is compatible and every triple
is incompatible~\cite[Example~3 and Sec.~VI]{AK}. The parent with at most one
reported label also occurs in the lossy binary construction of Bene and
V\'ertesi~\cite[Sec.~IV]{BV}. Zhang and Chen use the distinct-ray argument
to determine a parent for binary measurements from a rank-one
informationally complete POVM~\cite{ZC}.

\section{Quantitative stability}
For \(0\le\varepsilon\le1\), depolarization acts on an effect by
\[
   \D_\varepsilon(E)=(1-\varepsilon)E+\varepsilon\tau(E)\Id,
   \qquad \tau(E)=\tfrac12\tr(E).
\]
This unital positive map preserves every compatible set by acting on its
parent. To keep the nonfaces incompatible we use the overlap witness of
Skrzypczyk and Cavalcanti~\cite[Secs.~IV--V and Eq.~(12)]{SC}.
Their noise condition specializes to ours with detection efficiency
\(c\), visibility \(1-\varepsilon\), and dimension two. The facet labels
let us apply it to every nonface of a single measurement family.

\begin{theorem}[Depolarizing noise]\label{thm:noise}
Let \(\K\) have \(r\ge2\) facets. For the signed measurements
in \eqref{eq:facetchoice}, let
\[
 \mu=\max_{\substack{j\ne k\\s,t\in\{+,-\}}}
                 |\langle p_{j,s},p_{k,t}\rangle|,\qquad
 \delta=1-\mu>0 .
\]
Then \((\D_\varepsilon M_v)_{v\in V}\) has compatibility complex \(\K\)
whenever \(0\le\varepsilon<2\delta/r\). The bases can be chosen so that
every \(0\le\varepsilon<1/(2r^2)\) is allowed.
\end{theorem}
\begin{proof}
Fix a nonface \(S\). Every label occurs in \(S\), so choose \(v_j\in S\)
with \(j\in A_{v_j}\) for each \(j\). For any compatible family \(N\) on
\(S\) with these outcome alphabets, we claim
\begin{equation}\label{eq:witness}
 \mathcal W_S(N):=
 \sum_{j=1}^{r}\sum_{s=\pm}\tau\big[(P_{j,s}-\mu\Id)\,N_{v_j}(j,s)\big]
 \ \le\ \delta .
\end{equation}
On a fixed outcome string of a parent, let \(J\) contain the labels \(j\)
for which coordinate \(v_j\) reports \((j,s_j)\) for some sign \(s_j\).
For nonempty \(J\), the Gram matrix of the selected unit vectors has
absolute row sums at most \(1+(|J|-1)\mu\). Its largest eigenvalue is
also the largest eigenvalue of \(\sum_{j\in J}P_{j,s_j}\). Consequently
\[
   \sum_{j\in J}(P_{j,s_j}-\mu\Id)\preceq\delta\Id .
\]
The bound holds for empty \(J\) as well. Taking the normalized trace
against the positive parent effect on each string and summing proves
\eqref{eq:witness}. A coordinate selected for several labels can report
only one of them, so repeated choices of \(v_j\) cause no change.

For a rank-one qubit projector \(P\), use \(P^2=P\) and \(\tau(P)=1/2\) to get
\[
 \tau\big[(P-\mu\Id)\D_\varepsilon(cP)\big]
 =\frac c2\big[(1-\varepsilon)(1-\mu)+\varepsilon(1/2-\mu)\big]
 =\frac c2(\delta-\varepsilon/2).
\]
Summing over the two signs for every label gives
\begin{equation}\label{eq:gap}
     \mathcal W_S(\D_\varepsilon M)-\delta
       = r c(\delta-\varepsilon/2)-\delta
       = c(\delta-r\varepsilon/2).
\end{equation}
A positive gap excludes a joint measurement. This proves the first claim
for every nonface at once.

For the second claim, write a spin projector with Bloch direction
\(u\in\mathbb S^2\) as \(P(u)=(\Id+u\cdot\boldsymbol\sigma)/2\), where
\(\boldsymbol\sigma\) denotes the three Pauli matrices. Choose \(r\) axes
so that \(|u_j\cdot u_k|\le1-1/r\) for \(j\ne k\). Such a choice is
possible successively. For a fixed \(u_j\), the region
\(|u_j\cdot u|>1-1/r\) occupies \(1/r\) of the sphere, since \(u_j\cdot u\)
is uniform on \([-1,1]\) for uniform \(u\). The union of at most \(r-1\)
excluded regions has normalized area at most \((r-1)/r<1\), so another
axis can be chosen. As
\(\tr[P(u)P(v)]=(1+u\cdot v)/2\), both choices of sign give
\[
  \delta\ \ge\ 1-\sqrt{1-\frac1{2r}}\ \ge\ \frac1{4r}.
\]
\end{proof}

For an explicit signed realization with real effects, take the \(r\)
unoriented Bloch axes at angles \(j\pi/r\), \(0\le j<r\), in the \(xz\)-plane.
Then \(\delta=1-\cos(\pi/(2r))\), giving the sufficient noise threshold
\(\varepsilon<2[1-\cos(\pi/(2r))]/r\).
For \(r=2\), take two singleton facets and the orthogonal Pauli axes
\(X,Z\). The theorem guarantees incompatibility for \(\varepsilon<1-1/\sqrt2\).
At equality,
\[
        G(s,t)=\tfrac14\big[\Id+(sX+tZ)/\sqrt2\big],
        \qquad s,t\in\{-1,1\},
\]
is a parent. Mixing it with the uniform parent covers larger noise
fractions, giving the exact two-axis threshold~\cite{Busch}.

The merged realization has a positive margin as well. Set
\[
 \rho=e_1e_1^*,\qquad Q=\sum_j\tr(\rho P_j)>r-1,\qquad
 L=Q+\frac r2(\delta^{-1}-1)>0,
\]
where \(\delta\in(0,1]\) is the separation of the rational bases in
\eqref{eq:rational}, as defined in Theorem~\ref{thm:noise}.
For every nonface choose one vertex \(v_j\) reporting each positive label
\(j\). On any jointly measurable family \(N\) of these observables,
\begin{equation}\label{eq:merged-witness}
   \sum_j \tr\left[\left(\rho-\frac{\Id-P_j}{\delta}\right)
                         N_{v_j}(j)\right]\le 1 .
\end{equation}
If a parent string reports the selected labels \(J\ne\varnothing\),
the Gram bound gives \(\sum_{j\in J}(\Id-P_j)\succeq (|J|-1)\delta\Id\).
Combining this with \(\rho\preceq\Id\) gives
\[
 \sum_{j\in J}\left(\rho-\frac{\Id-P_j}{\delta}\right)
 \preceq |J|\rho-(|J|-1)\Id\preceq\rho .
\]
For \(J=\varnothing\) the sum is zero. Tracing against the positive
parent effects and summing gives \eqref{eq:merged-witness}.
After depolarization the left side equals \(c(Q-\varepsilon L)\).
Thus the merged pattern survives whenever
\begin{equation}\label{eq:mergedthreshold}
              0\le\varepsilon<(Q-r+1)/L .
\end{equation}

\begin{corollary}[Independent perturbations]\label{cor:perturb}
Let \(\K\) have \(r\ge2\) facets and let \(M\) be either realization.
Omit zero effects. Put \(\pi_v(x)=\tau(M_v(x))>0\) and \(f=\max_j|F_j|\). Suppose POVMs \(N_v(x)=\D_\varepsilon(M_v(x))+H_v(x)\)
on the same alphabets satisfy
\(\|H_v(x)\|_{\op}\le h\,\pi_v(x)\) for every \(v,x\). Their compatibility
complex is \(\K\) provided \(h\ge0\), \(0<\varepsilon\le1\), and
\begin{align}
 h&<\min\{\varepsilon/f,\ 2\delta/r-\varepsilon\}
       &&\text{for the signed realization},\label{eq:radius}\\
 h&<\min\left\{\varepsilon/f,\
        \frac{2(Q-r+1-\varepsilon L)}{r(1+1/\delta)}\right\}
       &&\text{for the merged realization}.\label{eq:mergedradius}
\end{align}
\end{corollary}
\begin{proof}
For a face \(S\), take a parent \(G\) of \((M_v)_{v\in S}\) on the
product outcome space. A parent for the perturbed family is
\[
 (1-\varepsilon)G(\boldsymbol x)
 +\Big(\prod_{v\in S}\pi_v(x_v)\Big)
   \left(\varepsilon\Id+\sum_{v\in S}\frac{H_v(x_v)}{\pi_v(x_v)}\right).
\]
The bracket is positive because \(|S|h<\varepsilon\). Since
\(\sum_x H_v(x)=0\), the second term contributes exactly
\(\varepsilon \pi_v(x_v)\Id+H_v(x_v)\) to marginal \(v\). For a nonface in the signed
realization, each selected outcome has \(\pi_{v_j}(j,s)=c/2\), and
\(\tau(|P_{j,s}-\mu\Id|)=1/2\). Hence the change in \eqref{eq:witness}
is at most \(rch/2\), smaller than the gap in \eqref{eq:gap}.
For the merged realization, each selected label again has \(\pi_{v_j}(j)=c/2\).
The trace norm of \(\rho-(\Id-P_j)/\delta\) is at most \(1+1/\delta\).
Thus \eqref{eq:mergedradius} bounds the change in \eqref{eq:merged-witness}
strictly below its gap \(c(Q-r+1-\varepsilon L)\).
\end{proof}

For positive \(\varepsilon\) below the corresponding noise threshold, the
strict bounds give an open neighborhood with pattern \(\K\) in the affine space
of Hermitian families on the fixed alphabets with \(\sum_x E_v(x)=\Id\)
for each \(v\). The lower bound \((\varepsilon-h)\pi_v(x)\Id\) makes all
effects there positive definite. For a simplex, the parent formula with
\(S=V\) and \(\varepsilon=1\) works near strictly positive scalar POVMs when
\(|V|h<1\).

\emph{Distinct settings.} For settings with at least two outcomes, coincidence under a fixed outcome
permutation is a proper affine condition: change two effects of one setting by \(+t\Id\) and \(-t\Id\) while holding the other setting fixed. Normalization is preserved and at most one value of \(t\) solves a specified coincidence. Each such affine subspace has empty interior, and a finite union of them has empty interior. The finitely many such conditions
can be avoided within the open neighborhood. In the real construction with
a rational noise fraction, rational symmetric perturbations suffice.
Vertices in every facet can use distinct biased binary coins within the
bounds for \(r\ge2\). Appending or forgetting independent coins preserves
compatibility.

\section{Why the outcome alphabet must grow}\label{sec:count}
A \emph{pairwise compatibility graph} has an edge \(ij\) when \(M_i\)
and \(M_j\) have a common parent. Grinko and Uola characterize compatibility
for families of unbiased binary qubit measurements~\cite{GU}. We count these
graphs at fixed dimension and outcome cap, allowing arbitrary bias.
One-outcome POVMs give only complete graphs, so we take \(q\ge2\).

\begin{theorem}[Fixed outcome cap]\label{thm:count}
Fix a Hilbert space dimension \(d\ge1\) and an outcome cap \(q\ge2\).
Let \(g_{d,q}(n)\) count the labelled pairwise compatibility graphs
realized by \(n\) POVMs with at most \(q\) outcomes each. With
\(m=d^2(q-1)\), there is a constant \(C_{d,q}\) such that, for all
sufficiently large \(n\),
\begin{equation}\label{eq:count}
       g_{d,q}(n)\le(C_{d,q}n)^{mn}.
\end{equation}
Consequently the proportion of labelled graphs so realizable tends to zero.
Thus for every fixed \(d\) and \(q\), some finite complexes have no realization
with \(q\) or fewer outcomes per measurement in dimension \(d\).
\end{theorem}
\begin{proof}
Pad each alphabet to \(q\) outcomes with zero effects. A point \(x\in\R^m\)
specifies the first \(q-1\) Hermitian matrices, and normalization gives the
last. Let \(R(x,y)\) hold when the two lists are POVMs with a common parent.
A parent consists of \(q^2\) positive semidefinite matrices \(G_{ab}\) whose
row and column sums equal the prescribed effects, enforcing their positivity.
These conditions are polynomial equalities and inequalities, using principal
minors for positivity. By Tarski--Seidenberg~\cite{Seidenberg}, their projection
\(R\) onto \((x,y)\) is semialgebraic. A fixed Boolean function of \(b\)
polynomial signs decides \(R\), with the degrees bounded by \(D\) and
\(b,D\ge1\) depending only on \(d,q\).

The graph depends on the signs of \(s=b\binom n2\) polynomials in \(N=mn\)
variables. As in the argument cited by Sauermann~\cite[App.~A.1]{Sauermann},
choose a point per realized pattern and a common \(\xi>0\) smaller than every
nonzero \(|p|\) at these points. The strict signs of the \(2s\) polynomials
\(p\pm\xi\) recover all original patterns. Warren's bound~\cite[Thm.~3]{Warren}
therefore gives, for \(s\ge N\),
\[
       g_{d,q}(n)\le (8eDs/N)^N
       = \left(\frac{4ebD(n-1)}{m}\right)^{mn},
\]
which implies \eqref{eq:count}.

There are \(2^{\binom n2}\) labelled graphs. Each is the pairwise graph
of a simplicial complex, for instance the complex consisting of its
edges, vertices and the empty face. Universality would require all these
graphs.
\end{proof}

For an explicit obstruction, let \(B_t\) be bipartite with parts \([t]\)
and \(2^{[t]}\), joining \(i\) to \(T\) exactly when \(i\in T\).
Fix the \(t\) left measurements. The signs of \(bt\) polynomials in the
\(m\) coordinates of a right measurement determine its neighbourhood.
The same bound allows \(O_{d,q}(t^m)\) possibilities. Since \(B_t\) requires
\(2^t\), it is impossible at any fixed \((d,q)\) for all sufficiently large
\(t\), including the binary qubit case \(d=q=2\).

The facet construction gives an outcome upper bound for each complex. The
counting argument shows that the worst-case outcome requirement diverges
with the number of vertices. Determining its growth rate remains open.

\paragraph{Acknowledgment.}
I used OpenAI Codex to develop and check proofs, run auxiliary
computations, and prepare this manuscript. The preprint is archived on
\href{https://doi.org/10.5281/zenodo.24528781}{Zenodo}.

\begin{thebibliography}{99}
\small
\setlength{\itemsep}{0pt}
\setlength{\parsep}{0pt}
\raggedright
\bibitem{KHF}
R.~Kunjwal, C.~Heunen and T.~Fritz,
Quantum realization of arbitrary joint measurability structures,
\href{https://doi.org/10.1103/PhysRevA.89.052126}{Phys.\ Rev.\ A
\textbf{89}, 052126 (2014)}.
\bibitem{AK}
N.~Andrejic and R.~Kunjwal,
Joint measurability structures realizable with qubit measurements:
Incompatibility via marginal surgery,
\href{https://doi.org/10.1103/PhysRevResearch.2.043147}{Phys.\ Rev.\ Research
\textbf{2}, 043147 (2020)}.
\bibitem{HDR}
R.~Kunjwal, \emph{Quantumness via Discrete Structures},
habilitation thesis, Aix-Marseille Universit\'e (2025), \href{https://arxiv.org/abs/2512.10063}{arXiv:2512.10063}.
\bibitem{SC}
P.~Skrzypczyk and D.~Cavalcanti, Loss-tolerant Einstein-Podolsky-Rosen
steering for arbitrary-dimensional states: Joint measurability and unbounded
violations under losses,
\href{https://doi.org/10.1103/PhysRevA.92.022354}{Phys.\ Rev.\ A
\textbf{92}, 022354 (2015)}.
\bibitem{BV}
E.~Bene and T.~V\'ertesi, Measurement incompatibility does not give rise
to Bell violation in general,
\href{https://doi.org/10.1088/1367-2630/aa9ca3}{New J.\ Phys.\
\textbf{20}, 013021 (2018)}.
\bibitem{ZC}
Y.-X.~Zhang and J.-L.~Chen, Measurement incompatibility and noiseless quantum codes,
\href{https://arxiv.org/abs/2610.37347}{arXiv:2610.37347 (2026)}.
\bibitem{Busch}
P.~Busch, Unsharp reality and joint measurements for spin observables,
\href{https://doi.org/10.1103/PhysRevD.33.2253}{Phys.\ Rev.\ D
\textbf{33}, 2253--2261 (1986)}.
\bibitem{GU}
D.~Grinko and R.~Uola, Compatibility of binary qubit measurements,
\href{https://doi.org/10.1103/vv1h-5mf9}{Phys.\ Rev.\ Lett.\
\textbf{135}, 200201 (2025)}.
\bibitem{Seidenberg}
A.~Seidenberg, A new decision method for elementary algebra,
\href{https://doi.org/10.2307/1969640}{Ann.\ of Math.\ (2)
\textbf{60}, 365--374 (1954)}.
\bibitem{Sauermann}
L.~Sauermann, On the speed of algebraically defined graph classes,
\href{https://doi.org/10.1016/j.aim.2021.107593}{Adv.\ Math.\
\textbf{380}, 107593 (2021)}.
\bibitem{Warren}
H.~E.~Warren, Lower bounds for approximation by nonlinear manifolds,
\href{https://doi.org/10.1090/S0002-9947-1968-0226281-1}{Trans.\ Amer.\
Math.\ Soc.\ \textbf{133}, 167--178 (1968)}.
\end{thebibliography}
\end{document}
